\chapter{Solutions to Chapter 2: Potential Outcomes}
\label{sol:chapter02}

\begin{exercisesolution}{A perfect doctor}{selection by treatment effect 下 observed difference in means 的偏差}
本题沿用 \cref{hw::perfect-doctor} 的设定。令
\[
X=Y(0)\sim \textsc{N}(0,1),\qquad \tau=-0.5+X,\qquad Y(1)=Y(0)+\tau=2X-0.5.
\]
perfect doctor 的 assignment rule 是
\[
Z=\mathbb{I}(\tau\geq 0)=\mathbb{I}(X\geq 0.5).
\]
因此 treatment group 正是 $X\geq 0.5$ 的 units，control group 正是 $X<0.5$ 的 units。Observed outcome 满足
\[
Y=
\begin{cases}
Y(1)=2X-0.5,& X\geq 0.5,\\
Y(0)=X,& X<0.5.
\end{cases}
\]
于是
\[
E(Y\mid Z=1)-E(Y\mid Z=0)
=
E(2X-0.5\mid X\geq 0.5)-E(X\mid X<0.5).
\]
对 standard Normal random variable，truncated mean 公式给出
\[
E(X\mid X\geq a)=\frac{\phi(a)}{1-\Phi(a)},\qquad
E(X\mid X<a)=-\frac{\phi(a)}{\Phi(a)}.
\]
取 $a=0.5$，得到
\[
E(Y\mid Z=1)-E(Y\mid Z=0)
=2\frac{\phi(0.5)}{1-\Phi(0.5)}-0.5+\frac{\phi(0.5)}{\Phi(0.5)}.
\]
用 $\phi(0.5)\approx 0.3521$、$\Phi(0.5)\approx 0.6915$，数值上
\[
E(Y\mid Z=1)-E(Y\mid Z=0)
\approx 2\frac{0.3521}{0.3085}-0.5+\frac{0.3521}{0.6915}
\approx 2.292 .
\]
这个值为正，而且很大。

但是 superpopulation average causal effect 是
\[
E\{\tau\}=E(-0.5+X)=-0.5.
\]
所以 observed difference in means 不但不等于 average causal effect，方向甚至相反。根本原因是 $Z$ 由 $\tau$ 本身决定，因而 treatment assignment mechanism 与 potential outcomes 强烈相关。换言之，treated units 是那些 control potential outcome 较高且 treatment effect 较大的 units；control units 是其余 units。比较 $E(Y\mid Z=1)$ 与 $E(Y\mid Z=0)$ 同时混合了 causal effect 与 selection bias。

\emph{Remark / 用途：} 这个练习是 \cref{eq::yobsdef,eq::fundamentalbridge} 的第一条重要警告。Observed outcome 总能写成 potential outcomes 与 assignment 的函数，但只看 observed outcome 的两个条件均值并不会自动识别 causal effect。后面 randomized experiments 的关键作用，正是用已知 assignment mechanism 切断这种 selection by potential outcomes。
\end{exercisesolution}

\begin{exercisesolution}{Nonlinear causal estimands}{median estimands 的三个例子与解释}
本题比较两个 nonlinear causal estimands：
\[
\delta_1=\text{median}\{Y_i(1)\}_{i=1}^n-\text{median}\{Y_i(0)\}_{i=1}^n,
\]
以及
\[
\delta_2=\text{median}\{Y_i(1)-Y_i(0)\}_{i=1}^n.
\]
为避免偶数样本的 median 约定影响例子，下面都取 $n=3$。

首先给出 $\delta_1=\delta_2$ 的例子。令
\[
(Y_i(0))_{i=1}^3=(0,1,2),\qquad (Y_i(1))_{i=1}^3=(1,2,3).
\]
则 individual effects 为 $(1,1,1)$，所以
\[
\delta_1=2-1=1,\qquad \delta_2=1.
\]
这里 treatment 对所有 units 都是 constant additive shift，因此两个 estimands 一致。

其次给出 $\delta_1<\delta_2$ 的例子。令
\[
(Y_i(0))_{i=1}^3=(0,10,11),\qquad (Y_i(1))_{i=1}^3=(1,2,12).
\]
则
\[
\text{median}\{Y_i(0)\}=10,\qquad \text{median}\{Y_i(1)\}=2,
\]
所以 $\delta_1=2-10=-8$。Individual effects 为
\[
(Y_i(1)-Y_i(0))_{i=1}^3=(1,-8,1),
\]
其 median 是 $1$，即 $\delta_2=1$。这个例子给出的是 $\delta_1<\delta_2$。

最后给出 $\delta_1>\delta_2$ 的例子。令
\[
(Y_i(0))_{i=1}^3=(0,1,12),\qquad (Y_i(1))_{i=1}^3=(-10,2,3).
\]
则
\[
\delta_1=2-1=1,
\]
而 individual effects 为
\[
(-10,1,-9),
\]
所以 $\delta_2=-9$。因此 $\delta_1>\delta_2$。

这些例子说明，median 与差分运算一般不能交换：
\[
\text{median}\{Y_i(1)\}-\text{median}\{Y_i(0)\}
\neq
\text{median}\{Y_i(1)-Y_i(0)\}.
\]
Average causal effect 是 linear estimand，所以差分和平均可以交换；median 是 nonlinear functional，因此没有这种代数性质。

哪一个 estimand 更合理，取决于 scientific question。若问题是“treatment 如何移动 outcome distribution 的中心位置”，尤其当两个 potential outcome distributions 可以看作同一人群在两个 regimes 下的 marginal distributions，那么 $\delta_1$ 很自然。例如，公共政策可能关心 treatment 是否提高总体收入分布的 median，因为 median income 是一个直接面向 population distribution 的指标。

若问题是“典型 unit 的 individual effect 有多大”，则 $\delta_2$ 更贴近问题本身。例如，个体化医疗中我们可能关心一个随机 patient 的 treatment benefit 的中位数：超过一半 patients 的 individual benefit 是否为正？这时 $\delta_2$ 直接描述 individual causal effects 的 distribution，而 $\delta_1$ 可能只反映两个 marginal distributions 的位置差。

不过，$\delta_2$ 有一个更强的概念要求：它依赖 individual-level joint schedule $(Y_i(1),Y_i(0))$，而 not only marginal distributions of $Y(1)$ and $Y(0)$。在实际数据中，两个 potential outcomes 不会同时观测到，因而识别 $\delta_2$ 往往比识别 $\delta_1$ 更难。这个区别会在后面对 causal estimands、randomization inference 和 partial identification 的讨论中反复出现。

\emph{Remark / 用途：} 这道题训练的是“先说清 estimand，再讨论 estimator”的习惯。很多应用争论并不是算法争论，而是目标量不同：population median shift、median individual benefit、quantile treatment effect 和 average treatment effect 回答的是不同问题。
\end{exercisesolution}

\begin{exercisesolution}{Average and individual effects}{positive average effect 但多数 units 不受益}
取 $n=3$，令
\[
(Y_i(0))_{i=1}^3=(0,0,0),\qquad
(Y_i(1))_{i=1}^3=(-1,-1,5).
\]
则 individual treatment effects 为
\[
(\tau_i)_{i=1}^3=(-1,-1,5).
\]
Average causal effect 等于
\[
\tau=\frac{-1-1+5}{3}=1>0.
\]
但是满足 $Y_i(1)>Y_i(0)$ 的 units 只有第 $3$ 个，比例为
\[
\frac{1}{3}<0.5.
\]
因此 treatment 的 average effect 为正，但少于一半的 units 从 treatment 中受益。

这个例子并不矛盾。Average causal effect 是 magnitude-weighted 的平均：一个非常大的正 effect 可以抵消多个较小的负 effects。相反，“受益者比例”是 sign-based estimand：
\[
n^{-1}\sum_{i=1}^n \mathbb{I}\{Y_i(1)>Y_i(0)\}.
\]
它只关心 effect 的方向，不关心大小。因此两者可以给出完全不同的 substantive message。

在 policy decision 中，这个区别很重要。如果 treatment 是低成本、低风险的政策，positive average effect 可能已经足够支持推广；如果 treatment 有严重副作用或公平性约束，那么“多数人是否受益”也许同样关键。医学中的个体化治疗、教育干预中的异质性响应、以及福利政策中的分配效应，都会要求同时报告 average effect 和 effect heterogeneity。

\emph{Remark / 用途：} 本题是 heterogeneity 的最小例子。它提醒我们，\cref{eq::teh,eq::teh2} 中的 $\tau_i$ 不是技术细节，而是 causal inference 的核心对象之一。Average effect 总结了总体收益，但不能替代 individual effects distribution。
\end{exercisesolution}
